\documentclass[10pt]{article}

\usepackage[top=0.55in, bottom=0.65in, left=0.55in, right=0.55in]{geometry}
\usepackage{amsmath, amssymb}
\usepackage[T1]{fontenc}
\usepackage[utf8]{inputenc}
\usepackage{lmodern}
\usepackage{enumitem}
\usepackage{multicol}
\usepackage{hyperref}

\setlength{\parindent}{0pt}
\setlength{\parskip}{0.35em}
\setlength{\columnsep}{0.22in}
\setlist[itemize]{leftmargin=*, topsep=2pt, itemsep=1pt, parsep=0pt}
\setlist[enumerate]{leftmargin=*, topsep=2pt, itemsep=1pt, parsep=0pt}

\newcommand{\KL}{\operatorname{KL}}
\newcommand{\clip}{\operatorname{clip}}
\newcommand{\InfoNCE}{\operatorname{InfoNCE}}
\newcommand{\softmax}{\operatorname{softmax}}
\newcommand{\E}{\mathbb{E}}
\newenvironment{symdefs}{\begin{itemize}[leftmargin=1.1em, topsep=1pt, itemsep=0pt, parsep=0pt]}{\end{itemize}}
\newcommand{\sym}[2]{\item \textbf{#1}: #2}

\begin{document}
\small

\begin{center}
{\Large \textbf{6.8610 Quiz 2 Study Note}}\\
Spring 2026\\
Aggregated from Lectures 8--15 in \texttt{6.8610/.quiz2}\\
Lightly cross-checked against \texttt{6.8610/.quiz2/NLP\_notes.pdf}
\end{center}

\begin{multicols*}{2}
\raggedcolumns

\section*{Cheatsheet Recommendations}
\begin{itemize}
\item \textbf{Use Side A for post-training.} Write the staged pipeline \texttt{pretrain -> mid-train -> SFT -> preference learning -> RL}, then the core equations for Bradley-Terry / DPO, the reward-maximization objective, and the PPO / GRPO update rules.
\item \textbf{Use Side B for retrieval and systems.} Put the retrieval tradeoff table \texttt{cross-encoder vs dual-encoder vs late interaction}, the InfoNCE loss, the ColBERT MaxSim score, the RAG factorization, and a tiny ReAct / CodeAct summary.
\item \textbf{Memorize comparisons, not paper names.} Expect ``why is method A better than B here?'' more than ``what year was paper X?''.
\item \textbf{High-yield formulas to handwrite:}
\[
\mathcal{L}_{\mathrm{DPO}},
\qquad
J(\theta)=\E[r]-\lambda \KL,
\qquad
\nabla_\theta \E[r]=\E[\nabla_\theta \log p_\theta \, r],
\]
\[
\mathcal{L}_{\InfoNCE},
\qquad
S_{\mathrm{ColBERT}}(q,d)=\sum_i \max_j E_{q_i}^{\top} E_{d_j},
\]
\[
p(a \mid q)=\sum_d p(d \mid q)p(a \mid q,d).
\]
\item \textbf{High-yield contrasts to handwrite:} SFT vs ranking losses vs RL; cross-encoder vs dual-encoder vs late interaction; closed-book QA vs RAG vs agents; byte-level vs BPE / WordPiece / Unigram.
\item \textbf{Always reserve space for failure modes.} This course repeatedly asks not just what works, but what breaks: SFT can teach style over truth, DPO can amplify noisy majority preferences, dense retrievers can fail OOD, RAG can memorize if retrieval is bad, agents suffer context rot, and tokenizer compression does not reliably predict downstream quality.
\item \textbf{If space is tight, drop examples before dropping objectives.} Keep one anchor example per lecture at most.
\item \textbf{Do not waste handwritten space on dates, author names, or leaderboard trivia.} Spend it on equations, assumptions, and ``what term does what?'' explanations.
\end{itemize}

\section*{Course Throughline}
\begin{itemize}
\item Lectures 8--10 explain how a pretrained next-token predictor becomes a useful assistant: better data staging, supervised demonstrations, preference data, and finally reward optimization with RL.
\item Lectures 11--12 explain how models access external knowledge: first by learning retrieval systems, then by combining retrieval with answer generation.
\item Lectures 13--14 move from ``model as chatbot'' to ``model as policy inside a system'' with tools, long-horizon interaction, prompt optimization, and modular AI programming abstractions.
\item Lecture 15 steps back and asks a representation question: before any of the above works, what basic discrete units should text be broken into?
\item Repeated theme: the right \textbf{inductive bias} often beats naive scaling. This is true for post-training, retrieval structure, agent scaffolding, and tokenization.
\end{itemize}

\section*{Lecture 8: Post-Training}
\textbf{Staged recipe:}
\[
\text{pretraining}
\rightarrow
\text{mid-training}
\rightarrow
\text{SFT}
\rightarrow
\text{preference learning / RL}.
\]

\textbf{SFT objective:}
\[
\max_\theta \sum_{(x,y)\in \mathcal{D}} \log p_\theta(y \mid x),
\]
usually with the loss masked on user tokens so the model is supervised mainly on assistant/output tokens.

\textbf{What mid-training is doing:}
\begin{itemize}
\item Continue next-token prediction, but on smaller, more curated data: code, books, technical writing, research papers, Q\&A, long-context or agent-style traces.
\item Early training can use shorter contexts for compute efficiency; later stages can increase context length once the model already has basic language priors.
\item High-quality data is scarce, so it is more valuable \emph{after} broad pretraining than before.
\item Later batches shape the model more strongly, so you would rather the model forget early boilerplate than later high-value technical data.
\end{itemize}

\textbf{What SFT is doing:}
\begin{itemize}
\item It does \emph{not} replace pretraining. It nudges and extracts capabilities already latent in the pretrained model.
\item It teaches instruction following, chat format, refusal style, and response structure.
\item Chat templates matter: tokens for system/user/assistant roles are part of the learned behavior.
\item If a new environment or tool setup appears later, the hope is that natural-language priors make ``post-post-training'' relatively cheap.
\end{itemize}

\textbf{Why naive SFT is limited:}
\begin{itemize}
\item It has a labeling bottleneck: you need prompt/desired-answer pairs.
\item It often forces one ``gold'' answer even when many answers would have been fine.
\item It can teach the \textbf{style} of competence without the deeper process that produced the answer.
\item It can worsen hallucination by pushing the model to sound authoritative rather than calibrated.
\item ``Say I don't know'' is often hard to teach from demonstrations alone.
\end{itemize}

\textbf{Exam intuition:} post-training works partly because pretrained models already contain broad language and world priors; later stages are steering those priors, not starting from scratch.

\section*{Lecture 9: Ranking Objectives for Human Feedback}
\textbf{Setup:}
\[
x=\text{prompt/context},
\qquad
y_w=\text{winner},
\qquad
y_\ell=\text{loser}.
\]

\textbf{SLiC / margin-ranking idea:}
\[
m(x,y_w,y_\ell)
=
\Delta(y_w,y_\ell)
+ \log p_\theta(y_\ell \mid x)
- \log p_\theta(y_w \mid x),
\]
\[
\mathcal{L}_{\mathrm{SLiC}}
=
\sum_{x,y_w,y_\ell} \max\!\left(0,m(x,y_w,y_\ell)\right)
+ \lambda \KL(p_\theta \,\|\, p_0).
\]
\begin{itemize}
\item \(\Delta\): required margin; can depend on similarity or task heuristics.
\item \(p_0\): reference or base model, often the SFT model.
\end{itemize}

\textbf{Bradley-Terry preference model:}
\[
P(y_w \succ y_\ell \mid x)
=
\sigma\!\left(r(x,y_w)-r(x,y_\ell)\right).
\]

\textbf{Bradley-Terry negative log-likelihood:}
\[
\mathcal{L}_{\mathrm{BT}}
=
-\sum_{x,y_w,y_\ell}
\log \sigma\!\left(r_\theta(x,y_w)-r_\theta(x,y_\ell)\right).
\]

\textbf{DPO parameterization:}
\[
r_\theta(x,y)
=
\lambda \log \frac{p_\theta(y \mid x)}{p_0(y \mid x)}
+ \lambda \log Z(x).
\]

\textbf{DPO loss:}
\[
\mathcal{L}_{\mathrm{DPO}}
=
-\sum_{x,y_w,y_\ell}
\log \sigma\!\left(
\lambda \log \frac{p_\theta(y_w \mid x)}{p_0(y_w \mid x)}
-
\lambda \log \frac{p_\theta(y_\ell \mid x)}{p_0(y_\ell \mid x)}
\right).
\]

\textbf{Unlikelihood for unpaired positive/negative data:}
\[
\mathcal{L}_{\mathrm{UL}}
=
-\sum_{(x,y^+)} \log p_\theta(y^+ \mid x)
-\sum_{(x,y^-)} \log \left(1-p_\theta(y^- \mid x)\right).
\]

\begin{itemize}
\item Ranking losses learn directly from \textbf{annotated responses}, not from new model samples.
\item Pairwise preference is often easier for humans than writing gold answers from scratch.
\item Feedback is especially helpful when a model should abstain or say ``I don't know'' rather than confidently imitate a bad pattern.
\item KL-to-reference prevents the model from satisfying only relative winner/loser orderings while drifting into garbage elsewhere.
\item DPO does \textbf{not} require a separate reward model; it uses LM probabilities plus a reference model.
\item KTO is a later variant inspired by asymmetric gain/loss ideas; the lecture treated it as lower priority than DPO.
\item Caveat: a single reward or preference function is a simplification; real users do not all want the same thing.
\end{itemize}

\textbf{Exam intuition:} DPO is elegant because it turns preference fitting into a classification-style loss while implicitly targeting a reward-regularized optimum.

\section*{Lecture 10: Reinforcement Learning Objectives}
\textbf{Why RL at all?} With limited preference data, it may be easier to learn \textbf{what good behavior looks like} than to directly learn the full policy that produces it.

\textbf{Reward-regularized objective:}
\[
J(\theta)
=
\sum_x
\E_{y \sim p_\theta(\cdot \mid x)}
\left[r(x,y)\right]
-\lambda \KL(p_\theta \,\|\, p_0).
\]

\textbf{Reward model:}
\[
r_\phi(x,y)=w^{\top} h_\phi(x,y),
\]
typically ``LM backbone + scalar head'', fit with the same Bradley-Terry likelihood:
\[
\max_\phi
\sum_{x,y_w,y_\ell}
\log \sigma\!\left(r_\phi(x,y_w)-r_\phi(x,y_\ell)\right).
\]

\textbf{Vanilla policy-gradient identity:}
\[
\nabla_\theta
\E_{y \sim p_\theta(\cdot \mid x)}[r(x,y)]
=
\E_{y \sim p_\theta(\cdot \mid x)}
\left[
\nabla_\theta \log p_\theta(y \mid x)\, r(x,y)
\right].
\]

\textbf{Interpretation:} sample from the model, score the sample, then do reward-weighted self-fine-tuning. If the reward is \(0/1\), this becomes ``keep the good samples, discard the bad ones''.

\textbf{Folding KL into the reward (useful approximation):}
\[
r'(x,y)
=
r(x,y)-\log p_\theta(y \mid x)+\log p_0(y \mid x).
\]

\textbf{Value and advantage:}
\[
V(x)=\E_{y \sim p_\theta(\cdot \mid x)}[r(x,y)],
\qquad
A(x,y)=r(x,y)-V(x).
\]

\textbf{Importance ratio for PPO:}
\[
\rho(x,y)=\frac{p_\theta(y \mid x)}{p_{\theta_{\mathrm{old}}}(y \mid x)}.
\]

\textbf{Clipped ratio:}
\[
\bar{\rho}(x,y)=\clip(\rho(x,y),1-\epsilon,1+\epsilon).
\]

\textbf{PPO objective:}
\[
\mathcal{L}_{\mathrm{PPO}}
=
\E_{y \sim p_{\theta_{\mathrm{old}}}}
\left[
\min\!\left(
\rho(x,y)A(x,y),
\bar{\rho}(x,y)A(x,y)
\right)
\right].
\]

\textbf{GRPO:}
for a group of samples \(y_1,\dots,y_n\) from the same prompt,
\[
\mu=\operatorname{mean}(r_i),
\qquad
\sigma=\operatorname{std}(r_i),
\qquad
A_i=\frac{r_i-\mu}{\sigma},
\]
then plug \(A_i\) into the same PPO-style clipped objective.

\begin{itemize}
\item Vanilla PG is unbiased but high variance and on-policy.
\item Baselines do not change the expectation; they reduce variance.
\item PPO adds three practical fixes: advantage estimation, sample reuse via importance sampling, and clipping for stability.
\item GRPO removes the learned value model and uses within-group reward normalization instead; this has become common in reasoning-model training.
\item The professor's key story: learn the reward from scarce human labels, then use extra compute to optimize that reward on \textbf{new model outputs} with no additional human annotation.
\end{itemize}

\textbf{Exam intuition:} vanilla PG is the conceptual core; PPO is the practical workhorse; GRPO is the simpler modern variant when many samples per prompt are available.

\section*{Lecture 11: Neural Information Retrieval}
\textbf{Task:} given query \(q\), corpus \(D\) with \(N\) documents, return top-\(k\) relevant documents quickly.

\textbf{Basic retrieval metrics:}
\[
\operatorname{Success@}k
=
\mathbf{1}[R_k(q)\cap G(q)\neq \emptyset],
\]
\[
\operatorname{Precision@}k
=
\frac{|R_k(q)\cap G(q)|}{k},
\qquad
\operatorname{Recall@}k
=
\frac{|R_k(q)\cap G(q)|}{|G(q)|}.
\]
\begin{symdefs}
\sym{\(R_k(q)\)}{retrieved top-\(k\) set}
\sym{\(G(q)\)}{gold relevant set}
\end{symdefs}

\textbf{Key complexity lens:}
\begin{itemize}
\item Cross-encoder: high quality, but \(O(N)\) Transformer calls per query.
\item Dual-encoder: \(O(1)\) Transformer calls per query, but still \(O(N)\) dot products.
\item ANN / IVF: heuristic route toward roughly \(O(\sqrt{N})\) candidate scoring under balanced clustering assumptions.
\end{itemize}

\textbf{Dual-encoder scoring:}
\[
s(q,d)=f(q)^{\top}g(d).
\]

\textbf{InfoNCE contrastive loss:}
\[
\mathcal{L}_{\InfoNCE}
=
-\log
\frac{
\exp(f(q)^{\top}g(d^+)/\tau)
}{
\exp(f(q)^{\top}g(d^+)/\tau)
+\sum_i \exp(f(q)^{\top}g(d_i^-)/\tau)
}.
\]
\begin{symdefs}
\sym{\(d^+\)}{positive document}
\sym{\(d_i^-\)}{sampled negatives}
\sym{\(\tau\)}{temperature}
\end{symdefs}

\textbf{Late interaction / ColBERT score:}
\[
S(q,d)=\sum_i \max_j E_{q_i}^{\top} E_{d_j}.
\]
\begin{symdefs}
\sym{\(E_{q_i}\)}{embedding of query token \(i\)}
\sym{\(E_{d_j}\)}{embedding of document token \(j\)}
\end{symdefs}

\textbf{IVF sketch:}
\begin{itemize}
\item Cluster document vectors into about \(\sqrt{N}\) buckets.
\item At query time, compare to \(\sqrt{N}\) centroids, then score only docs in the nearest bucket(s), about \(O(\sqrt{N})\) docs if clusters are balanced.
\end{itemize}

\begin{itemize}
\item Cross-encoders are great rerankers, not good first-stage retrievers.
\item Hard negatives matter more than random negatives because random negatives quickly become too easy.
\item Single-vector dense retrieval often struggles most on \textbf{generalization / OOD robustness}, not just expressivity.
\item Late interaction is the ``missing middle'': independent precomputable encodings plus fine-grained token-level matching.
\item ColBERT storage blowup is handled by smaller vectors plus aggressive compression / quantization.
\item ANN still works for late interaction because a doc can only score highly if at least one of its token vectors gets close to a query token.
\end{itemize}

\textbf{Exam intuition:} the design move is repeated decomposition: across documents, then between query and document, then across tokens inside a document.

\section*{Lecture 12: Retrieval-Augmented Generation}
\textbf{Closed-book open-domain QA:}
answer from model parameters alone.

\textbf{Open-book open-domain QA:}
store knowledge explicitly in documents and learn to retrieve and use them.

\textbf{RAG factorization:}
\[
p(a \mid q)
=
\sum_{d \in \mathcal{D}}
p_\eta(d \mid q)\, p_\theta(a \mid q,d).
\]
In practice, this is approximated over top-\(K\) retrieved documents.

\textbf{Reader styles:}
\begin{itemize}
\item Extractive reader: predict an answer span from retrieved text.
\item Generative reader: decoder generates the answer conditioned on retrieved text; can paraphrase or combine evidence across documents.
\end{itemize}

\textbf{Why closed-book QA is not enough:}
\begin{itemize}
\item Parametric knowledge is static.
\item Provenance is opaque.
\item Private or fresh corpora may never appear in pretraining.
\item Using parameters to memorize huge long-tail fact sets may be inefficient.
\end{itemize}

\textbf{Retriever-supervision problem:}
\begin{itemize}
\item Often the dataset gives only \((q,a)\), not gold relevant documents.
\item Heuristics include: use an existing retriever, ask humans for relevance, or mark answer-containing documents as pseudo-positives.
\item But search relevance and answer usefulness are related, not identical.
\end{itemize}

\textbf{Why latent-document marginalization was attractive:}
\begin{itemize}
\item Better credit assignment than RL when the retriever found the right doc but the reader still failed.
\item But it is hard to initialize, can push the reader to memorize answers if top-\(K\) is junk early on, and makes index refresh expensive.
\end{itemize}

\textbf{Fusion-in-Decoder (FiD):}
\begin{itemize}
\item Encode each retrieved passage separately.
\item Let the decoder attend over all passage encodings together.
\item This buys long effective context without one giant quadratic joint encoder.
\end{itemize}

\textbf{Multi-hop / deep research:}
\begin{itemize}
\item Earlier benchmarks: multi-hop factoid QA.
\item Modern framing: ``deep research'' systems that iteratively search, browse, condense evidence, and answer.
\item Baleen-style systems combat combinatorial explosion via iterative retrieval plus evidence condensation / compaction.
\item A major lecture pivot: maybe retrieval is best viewed as a \textbf{tool} the LM learns to use in text space, rather than something always trained end-to-end with the generator.
\end{itemize}

\textbf{Exam intuition:} RAG is fundamentally about \textbf{credit assignment} and \textbf{where knowledge should live}: in weights, in retrieved context, or in a policy that knows how to query tools.

\section*{Lecture 13: Agents and Inference Scaling}
\textbf{Inference-time scaling idea:}
some tasks improve if the model is allowed to think longer, sample more trajectories, or use tools before committing to a final answer.

\textbf{Pass@1 vs pass@\(k\):}
\begin{itemize}
\item \textbf{pass@1}: expected accuracy from a single rollout.
\item \textbf{pass@\(k\)}: success if at least one of \(k\) rollouts solves the task.
\end{itemize}

\textbf{Agent loop abstraction:}
\[
s_t
\xrightarrow{\text{LM policy}}
a_t
\xrightarrow{\text{environment/tool}}
o_{t+1}
\rightarrow
s_{t+1}.
\]

\textbf{Tool use in token space:}
\begin{itemize}
\item Model emits special tool-call syntax.
\item Generation is interrupted.
\item External system executes the tool.
\item Observation is appended to context.
\item Generation resumes.
\end{itemize}

\textbf{ReAct vs CodeAct:}
\begin{itemize}
\item \textbf{ReAct}: interleave ``thought'' tokens with explicit actions; reasoning can be a no-op action that buys more inference-time compute.
\item \textbf{CodeAct}: take actions by writing code; much more general, but more dangerous and harder to control.
\end{itemize}

\textbf{What an agent is in this lecture:}
\begin{itemize}
\item Not a sharply defined object category.
\item A useful abstraction for systems that iteratively pursue goals inside environments.
\item Strong post-trained LMs can already look agentic if tools and observations are expressed in language.
\end{itemize}

\textbf{Long-horizon problems:}
\begin{itemize}
\item Large environments, especially codebases.
\item Persistent state changes.
\item Long contexts cause cost and context rot.
\item Context condensation / compaction keeps the active context small while preserving enough history to continue.
\end{itemize}

\textbf{Post-training for agent behavior:}
\begin{itemize}
\item Mid-train heavily on high-quality code and tool-use traces.
\item Train across many environments, not just one.
\item Use RL on full trajectories when rewards are sparse and long-horizon.
\item An agent must stay anchored to the \textbf{goal}, not just predict locally likely next tokens.
\end{itemize}

\textbf{Exam intuition:} agents generalize RAG. Instead of hard-coding all outside computation, we give the model some control over when and how to invoke tools.

\section*{Lecture 14: Language Model Programming}
\textbf{Main shift:}
\[
\text{ad-hoc prompt} \rightarrow \text{reliable AI system}.
\]

\begin{itemize}
\item A one-off chatbot interaction is not the same problem as engineering a maintainable, transparent, reliable system that solves the same kind of task over and over.
\item Symbolic code is useful because it gives reusable modules, explicit dataflow, and clearer contracts.
\item Natural language is still necessary because many requirements are easiest to specify in English rather than in labels alone.
\end{itemize}

\textbf{Central abstraction: natural-language signatures.}
\begin{itemize}
\item Example mental model: \texttt{email thread -> events, action items}.
\item The point is to separate \textbf{what} the module should do from \textbf{how} a particular model or inference pipeline currently does it.
\item This localizes ambiguity inside module boundaries instead of smearing it across one giant fragile prompt.
\end{itemize}

\textbf{Three ingredients of an AI program:}
\begin{itemize}
\item \textbf{Code / modular structure:} how modules compose.
\item \textbf{Natural-language specs:} task definitions and output contracts.
\item \textbf{Data / feedback:} edge cases, rewards, and examples of where the first two are insufficient.
\end{itemize}

\textbf{Compile the intent into optimized artifacts:}
\begin{itemize}
\item optimize prompts,
\item optimize weights with RL or fine-tuning,
\item optimize the pipeline / inference strategy itself.
\end{itemize}

\textbf{Prompt optimization themes:}
\begin{itemize}
\item Bootstrap useful trajectories and search for prompts that make them reusable.
\item MIPRO: smarter search over prompt variants instead of naive random search.
\item TextGrad / GEPA style methods: use an LM to reflect on failures, propose textual fixes, and iteratively evolve prompts.
\item Prompt optimization can be far more sample-efficient than policy-gradient RL when the model is already strong enough to reason about what went wrong.
\end{itemize}

\textbf{But prompt optimization is limited:}
\begin{itemize}
\item It depends on the model already being capable of following the improved instructions.
\item It cannot inject missing domain knowledge the model simply does not have.
\item Hence prompt optimization and RL / fine-tuning are complements, not substitutes.
\end{itemize}

\textbf{Programming abstraction message:}
\begin{itemize}
\item We should not hard-wire every model-specific prompting trick into application logic.
\item Inference strategies like ReAct should be abstracted so they can compose over arbitrary signatures and be learned / optimized too.
\item A good abstraction hides model-specific plumbing, just like PyTorch hides low-level tensor kernels.
\end{itemize}

\textbf{Exam intuition:} this lecture is about moving from ``prompt engineering hacks'' to a declarative view where AI systems are specified at a higher level and then compiled down into prompts, weights, and pipelines.

\section*{Lecture 15: Tokenization}
\textbf{Representation problem:} unlike images or audio, text does not come with a natural low-level numeric representation whose distances already mean something useful.

\textbf{Byte-level tokenization:}
\begin{itemize}
\item Finite vocabulary, so no OOV problem.
\item But sequences become long.
\item Byte values are arbitrary with respect to meaning.
\item Byte-level generation can drift into gibberish because bytes are not word-like units.
\end{itemize}

\textbf{Word-level tokenization:}
\begin{itemize}
\item Shorter sequences and more intuitive units.
\item But morphology, unseen words, multilingual variation, and even the notion of ``word'' itself make it brittle.
\end{itemize}

\textbf{Idealized tokenizer objective:}
\[
\min_{\Sigma,f}
\sum_i |f(x^{(i)})|
\qquad
\text{subject to } |\Sigma|=V.
\]
\begin{symdefs}
\sym{\(\Sigma\)}{token vocabulary}
\sym{\(f\)}{tokenization function}
\sym{\(V\)}{target vocabulary size}
\end{symdefs}

\textbf{Why this is not solved exactly:}
\begin{itemize}
\item A perfect tokenizer would balance compactness and expressivity.
\item But globally optimal tokenization under natural formulations is intractable, so practical tokenizers use greedy approximations.
\end{itemize}

\textbf{Byte-Pair Encoding (BPE):}
\begin{itemize}
\item Start from bytes / characters.
\item Repeatedly merge the most frequent adjacent pair.
\item Store the merge history \(M\).
\item At test time, apply those learned merges to new text.
\end{itemize}

\textbf{BPE lesson from the lecture:}
\begin{itemize}
\item BPE is a compromise: words are too coarse, bytes are too fine.
\item It often compresses token count well, but the learned tokens are not guaranteed to be linguistically satisfying.
\item Greedy longest-available segmentation tends to work well at test time.
\end{itemize}

\textbf{Pretokenization:}
\begin{itemize}
\item Hand-built guardrails often stop merges across spaces, punctuation, and digit chunks.
\item This is not elegant, but it is practical and important.
\item Real production tokenizers use heavy regex-based pretokenization.
\end{itemize}

\textbf{Other classics:}
\begin{itemize}
\item \textbf{WordPiece}: BPE-like but merge choice is likelihood / PMI flavored rather than raw count based.
\item \textbf{Unigram LM}: start with a large vocabulary and prune the least useful tokens.
\item \textbf{SentencePiece}: framework that supports BPE and Unigram, often useful when whitespace assumptions are undesirable.
\end{itemize}

\textbf{What the later ``advances'' section argued:}
\begin{itemize}
\item Compression alone does \textbf{not} predict downstream performance well.
\item Vocabulary size effects can be surprisingly weak over a normal range.
\item Pretokenization choices matter a lot.
\item No tokenizer clearly dominates all others.
\item Byte-level / tokenizer-free approaches are becoming more competitive again.
\item Special tokens for system/user/assistant roles are inserted during instruction tuning and are not split into smaller subwords.
\end{itemize}

\textbf{Exam intuition:} tokenization is a representation choice, not a solved preprocessing detail. The lecture's conclusion is basically ``we still do not know what the perfect tokenizer should optimize''.

\section*{Likely Quiz Habits}
\begin{itemize}
\item Be able to write each objective and explain what every term is doing in plain English.
\item Expect compare/contrast questions framed around sample efficiency, latency, credit assignment, or inductive bias.
\item When asked ``why does this help?'', answer with the bottleneck it addresses: better supervision, better reward credit, better retrieval structure, better tool use, or better representation.
\item For retrieval and agents, mention systems constraints alongside model quality. This part of the course is explicitly about quality \emph{and} latency / environment interaction.
\item For tokenization, do not overclaim. The lecture repeatedly emphasized uncertainty: good compression is not the same thing as a universally good tokenizer.
\end{itemize}

\end{multicols*}

\end{document}
